\documentclass[10pt,a4paper]{amsart}
\usepackage[margin=19mm]{geometry}
\usepackage{amsmath,amssymb,mathtools}
\usepackage[scr=rsfs,cal=euler]{mathalfa}
\usepackage{xfrac}
\usepackage[unicode,hidelinks,bookmarks=false]{hyperref}
\newcommand{\ie}{i.e.\ }
\newcommand{\cf}{cf.\ }
\newcommand{\resp}{respectively\ }
\newcommand{\Subsection}{\subsection}
\providecommand{\nobreakdash}{\nobreak\hbox{-}}
\setlength{\emergencystretch}{2em}
\multlinegap0pt
\usepackage{bm}
\usepackage{bbm}
\usepackage{dsfont}
\usepackage{xspace}
\usepackage{cancel}
\usepackage{mathtools}
\usepackage{amsthm}
\makeatletter
\@ifundefined{defn}{\newtheorem{defn}{Defn}}{}
\@ifundefined{Prop}{\newtheorem{Prop}[defn]{Proposition}}{}
\@ifundefined{claim}{\newtheorem*{claim}{Claim}}{}
\@ifundefined{fact}{\newtheorem{fact}[defn]{Fact}}{}
\@ifundefined{lem}{\newtheorem{lem}[defn]{Lemma}}{}
\makeatother
\newcommand{\BE}{\mathrm{BE}}
\newcommand{\C}{\mathbb{C}}
\title[Integration in finite terms and exponentially algebraic func]{Integration in finite terms and exponentially algebraic functions}
\author{R\'emi Jaoui, Jonathan Kirby}
\date{Source: arXiv:2510.26248v1 (CC BY 4.0)}
\begin{document}
\maketitle

\noindent\emph{Source and licence.} Integration in finite terms and exponentially algebraic functions. Version arXiv:2510.26248v1 (CC BY 4.0), distributed under the Creative Commons Attribution 4.0 International licence, as recorded in the arXiv licence metadata for this exact version and shown on the abstract page https://arxiv.org/abs/2510.26248v1. Full rights notice: licenses/arxiv-2510.26248-CC-BY-4.0.txt.\par\smallskip

\noindent\textit{Editorial scope.} Counted unit: the Prop whose source label is \texttt{o-minimal direction1} in the article (source0.tex lines 892--894, source label \texttt{o-minimal direction1}), delivered together with the article's own complete original proof of it (source0.tex lines 896--931, one \texttt{proof} environment of 36 lines). No mathematics is written, repaired or re-derived: statement and proof are the article's own text line for line. Required context retained uncounted: lemma-exponentialalgebra (lines 590--595); fact-Wilkie-complete (lines 787--802); lemma-generic-points (lines 820--824); prop-from relation to functions (lines 842--845). Every definition, display and label of those lines is retained unchanged. Omitted from this package and not redistributed: 1--589, 596--786, 803--819, 825--841, 846--891, 895--895, 932--1914. Nothing inside a retained excerpt is omitted or altered: every retained body is delivered exactly as the article's lines are written, so no omission receipt of any kind accompanies this package. Citations: the retained excerpts carry no citation command at all, so no bibliography entry of the article is transcribed and no reference is redistributed. Reference adaptation: none was needed and none was made; every reference inside the delivered unit resolves inside it, so no unresolved reference or citation is produced. Printed slips kept literal: no printed word, symbol, hypothesis or number of the counted statement or of its proof was altered. Layout and class: the article's own class is \texttt{amsart} with packages including amsmath, amssymb, amsfonts, graphics, natbib, babel, tikz, inputenc, xy, fontenc, pifont, geometry, hyperref; the wrapper is the lane's pinned \texttt{amsart} editorial wrapper at 10pt on A4, which restates the theorem-like environments the retained text needs and adds the article's own macro definitions that the retained text needs (\texttt{\textbackslash BE}, \texttt{\textbackslash C}), with extra wrapper packages (bm, bbm, dsfont, xspace, cancel, mathtools). The delivered numbering is the unit's own. Assets: no retained span contains a figure environment, a graphics-inclusion command, a PDF-page inclusion, a table or a TikZ picture; no image file is copied into this package, the package declares no asset dependency and no image file is redistributed with it. Licence: version arXiv:2510.26248v1 of this article is distributed under the Creative Commons Attribution 4.0 International licence, as recorded in the arXiv licence metadata for this exact version and shown on the abstract page https://arxiv.org/abs/2510.26248v1. A full rights notice is delivered at licenses/arxiv-2510.26248-CC-BY-4.0.txt. Characters: every retained excerpt is pure ASCII, so the delivered engine, XeLaTeX with its default Latin Modern text font, reports no missing character.
\begin{lem}\label{lemma-exponentialalgebra} 
Let $L/k$ be an extension of blurred exponential fields with the same field of constants. The space $\mathrm{Der}_{\BE}(L/k)$ of $\BE$-derivations is an $L$-vector space. A lower bound for its dimension is given by:
$$ \delta(L/k) \leq \mathrm{ldim}_L(\mathrm{Der}_{\BE}(L/k)).$$
Furthermore, if $k$ is self-sufficient in $L$ in the sense that $\delta(M/k) \geq 0$ for all intermediate extension contained in $L$ then 
$$ \delta(L/k) = 0 \text{ if and only if } \mathrm{Der}_{\BE}(L/k) = \lbrace 0 \rbrace.$$
\end{lem}

\begin{fact} \label{fact-Wilkie-complete}
Let $k_0$ be a countable algebraically closed subfield of parameters.
The following properties hold:
\begin{itemize} 
\item[(a)] $\mathrm{hcl}_{k_0}(-): \mathcal P(\mathbb{C}) \rightarrow \mathcal P(\mathbb{C})$ is a pregeometry on $\mathbb{C}$. Furthermore, $a_1,\ldots,a_n \in \mathbb{C}$ are $\mathrm{hcl}$-independent if and only if for any $k_0$-definable holomorphic function $F$ with $(a_1,\ldots,a_n) \in \mathrm{dom}(F)$ then 
$$ F(a_1,\ldots,a_n) = 0 \Rightarrow F = 0.$$

\item[(b)] If $g$ is a $k_0$-definable holomorphic function of the variables $z = (z_1,\ldots, z_n)$ and $ b = (b_1,\ldots,b_n) \in \mathrm{dom}(g)$ are $\mathrm{hcl}$-independent then $g$ can be $m$-implicitly defined from $\mathcal A$ in a neighborhood $W$ of $b$ in the following sense: 

\begin{quote} 
there is $m \geq 1$, holomorphic functions $g_1 = g, g_2,\ldots,g_m$ on $W$ and $F_1,\ldots,F_m \in \mathcal A_{n + m} = \mathbb{C}[z,e^{z},w,e^{w}]$ such that for all $z \in W$ 
$$ F_i\Big(z,e^{z}, g_1(z), \ldots, g_m(z), e^{g_1(z)}, \ldots, e^{g_m(z)} \Big) = 0 $$
for $i = 1,\ldots, m$ and $\mathrm{det}\Big((\frac{\partial F_i} {\partial w_j})_{i,j \leq m}(\overline{z}, g_1(\overline{z}),\ldots, g_m(\overline{z})) \Big)\neq 0$.
\end{quote}
\end{itemize}
\end{fact}

\begin{lem} \label{lemma-generic-points}

Let $f_1,\ldots,f_n$ be holomorphic functions defined on some complex domain $U$ which are $\mathrm{hcl}$-independent over $\mathbb{C}$ and let $k_0$ be a countable subfield of $\mathbb{C}$. Then the analytic curve $(f_1,\ldots,f_n)(U)$ passes through an $\mathrm{hcl}$-generic point, that is, there exists $a \in U$ such that the values
$f_1(a),\ldots,f_n(a)$ are $\mathrm{hcl}$-independent.
\end{lem} 

\begin{Prop} \label{prop-from relation to functions}
Let $f_1,\ldots,f_n,f$ be holomorphic functions. Assume that $f_1,\ldots,f_n$ are $\mathrm{hcl}$-independent over $\mathbb{C}$ and that $f_1,\ldots,f_n,f$ are not. Then there exists a subdomain $V$ of $U$, $W$ a connected open subset of $\mathbb{C}^n$ and a definable holomorphic function $G$ on $W$ such that
$$ f_{\mid V} = G \circ (f_{1 \mid V},\ldots,f_{n \mid V}) $$
\end{Prop}

\begin{Prop}\label{o-minimal direction1}
If holomorphic functions of one variable  $f_1,\ldots,f_{n}$ are $\mathrm{hcl}$-dependent over $\mathbb{C}$ then  they are also  exponentially algebraically dependent over $\mathbb{C}$.
\end{Prop} 

\begin{proof} 
Work by induction and assume that the statement holds for some $n$ and consider holomorphic functions $f_1,\ldots,f_{n+1}$ which are $\mathrm{hcl}$-dependent over $\mathbb{C}$. If already $f_1,\ldots,f_n$ are $\mathrm{hcl}$-dependent over $\mathbb{C}$, then the induction hypothesis implies that $f_1,\ldots,f_{n}$ are exponentially algebraically dependent over $\mathbb{C}$. So we may assume without loss of generality that $f_1,\ldots,f_n$ are both $\mathrm{hcl}$-independent over $\mathbb{C}$ and $\mathrm{ecl}$-independent over $\mathbb{C}$. 

Hence, Proposition \ref{prop-from relation to functions} applies and after replacing $U$ by a subdomain, we can write 
\begin{equation}\label{equality-proof-equivalence} f_{n+1} = G \circ (f_1,\ldots,f_n) 
\end{equation}
where $G$ is a definable holomorphic function defined on some $W \subset \mathbb{C}^n$ containing $(f_1,\ldots,f_n)(U)$. Fix $k_0$ a countable field of definition of $G$. By Lemma \ref{lemma-generic-points}, we can find $a \in U$ such that the values $b_1 = f_1(a), ... , b_n = f_n(a)$ are $\mathrm{hcl}$-independent. Hence Fact \ref{fact-Wilkie-complete} (b) applies at the point $(b_1,\ldots,b_n)$ and we can find $G = G_1,\ldots,G_m$ satisfying the conclusion of Fact \ref{fact-Wilkie-complete} (b). Up to restricting the domain $U$ again, the compositions
$$g_i = G_i \circ (f_1,\ldots, f_n) \text{ and } h_i = e^{g_i}$$
are well-defined holomorphic functions on $U$ for $i =1,\ldots,m$. Set 
$$ k = \mathbb{C}(f_1,\ldots,f_n,e^{f_1},\ldots,e^{f_n})^{alg} \text{ and } L = k(g_1,\ldots,g_m, h_1,\ldots,h_m)^{alg}.$$

\begin{claim} 
$\delta(L/k) = 0.$
\end{claim} 


\begin{proof}[Proof of the claim]

By Lemma \ref{lemma-exponentialalgebra}, we need to show that $\mathrm{Der}_{\BE}(L/k)$ is the trivial vector space so take $D \in \mathrm{Der}_{\BE}(L/k)$.  Since $G_1,\ldots,G_m$ are implicitly defined from $\mathbb{C}[z,w,e^z,e^w]$ at the point $b$, it follows that we can find polynomials $$P_1,\ldots,P_m \in k[W_1,\ldots,W_m, E_1,\ldots,E_m]$$ such that $P_s(g_1,\ldots,g_m,h_1,\ldots,h_m) = 0$. Viewing the $P \in k[W_1,\ldots, W_m, E_1,\ldots, E_m]$ as an exponential polynomial given by:
$$ P \mapsto \overline{P} =  P(W_1,\ldots,W_n,e^{W_1},\ldots, e^{W_n})  $$
the partial derivative along the $W_i$ can be read as 
$$ \frac{\partial \overline{P}}{\partial W_i} = \overline{\frac{\partial P} {\partial W_i} + E_i \cdot \frac {\partial P} {\partial E_i}}.$$ 
Setting $\partial_i = \frac{\partial} {\partial W_i} + E_i \cdot \frac {\partial} {\partial E_i} \in \mathrm{Der}(k[W_1,\ldots,W_m,E_1,\ldots,E_m]/k)$ the condition on the nonvanishing of the Jacobian in Fact  \ref{fact-Wilkie-complete} (b) can be written as
\begin{equation}\label{equation-nonvanishingoftheJacobian}
 \mathrm{det}\Big( \partial_i(P_s) (g_1,\ldots,g_m,h_1,\ldots,h_m) \Big) \neq 0 \in L.
 \end{equation}
Now since $D$ is a $\BE$-derivation, we have that $D(h_i) = h_i D(g_i)$ so that differentiating $P_s(g_1,\ldots,g_m,h_1,\ldots,h_m) = 0$ we get 
\begin{eqnarray*}
0  & = &  \sum_{i = 1}^m \frac{\partial P_s}{\partial W_i}(g,h) D(g_i) + \sum_{i = 1}^m \frac{\partial P_s}{\partial E_i}(g,h) h_i D(g_i) \\
& = & \sum_{i = 1}^m \Big( \frac{\partial P_s}{\partial W_i}(g,h) + \frac{\partial P_s}{\partial E_i}(g,h) h_i \Big) D(g_i) \\
& = &  \sum_{i = 1}^m \partial_i(P_s)(g,h) D(g_i).
\end{eqnarray*}
Since this is true for any $s = 1,\ldots, m$, the vector $D(g_1),\ldots, D(g_m)$ lies in the kernel of the matrix $(\partial_i(P_s)(g,h))_{i,s}$ which is trivial by (\ref{equation-nonvanishingoftheJacobian}). It follows that $D(g_1) = \cdots = D(g_m) $ so that $D(h_1) = \ldots = D(h_m) = 0$ and $D = 0$ as required.
\end{proof}
To conclude the proof, note that since $f_1,\ldots,f_n$ are exponentially algebraically independent over $C$, $k$ is a self-sufficient subfield of $K$. Since Equation (\ref{equality-proof-equivalence}) implies that $f_{n+1} = g_1 \in L$ and we have shown that $\delta(L/k) = 0$, we conclude that $f_{n+1} \in \mathrm{ecl}^{K}(\C, f_1,\ldots,f_n)$. This means that $f_1,\ldots, f_{n+1}$ are exponentially algebraically dependent over $\mathbb{C}$.
\end{proof} 

\end{document}
